% A made-up paper for Kovrune's mathematics suite (eval/math/README.md). The flow, its results and its
% references are invented; it is written as an arXiv analysis-of-PDEs paper is: thmtools environments
% sharing one counter numbered within sections, main theorems lettered A and B, hypotheses tagged (H1)
% and (H2), and equations numbered within sections.
\documentclass[11pt]{article}
\usepackage{amsmath,amssymb,amsthm,thmtools}
\usepackage{hyperref}
\usepackage{cleveref}

\declaretheorem[name=Theorem, numberwithin=section]{theorem}
\declaretheorem[name=Lemma, sibling=theorem]{lemma}
\declaretheorem[name=Proposition, sibling=theorem]{proposition}
\declaretheorem[name=Conjecture, sibling=theorem]{conjecture}
\declaretheorem[name=Remark, style=remark, sibling=theorem]{remark}
\declaretheorem[name=Theorem]{main}
\renewcommand{\themain}{\Alph{main}}

\numberwithin{equation}{section}

\newcommand{\R}{\mathbb{R}}
\newcommand{\qc}{q_{*}}

\title{Decay and blow-up for the fractional Odum flow}
\author{Lucia Ferreira-Hahn}
\date{}

\begin{document}

\maketitle

\begin{abstract}
We study the fractional Odum flow $\partial_t u + (-\Delta)^s u + V u = |u|^{q-1} u$ on $\R^d$ with a nonnegative potential $V$. For $d \geq 3$ and $1/2 < s < 1$ we prove that small solutions are global and decay when $q$ exceeds the critical exponent $\qc = 1 + 4s/(d+2)$, and that every nonnegative nontrivial solution blows up in finite time when $1 < q < \qc$. The critical case $q = \qc$ remains open; we conjecture that solutions blow up there too.
\end{abstract}

\section{Introduction}\label{sec:intro}

Let $d \geq 3$ and $0 < s < 1$. The fractional Odum flow is the equation
\begin{equation}\label{eq:flow}
\partial_t u + (-\Delta)^s u + V u = |u|^{q-1} u, \qquad u(0) = u_0,
\end{equation}
for $u \colon [0, T) \times \R^d \to \R$, where $q > 1$ and $V \colon \R^d \to \R$ is a potential. The equation was introduced by Odum as a model of anomalous diffusion in a medium with absorption. Its behaviour is decided by the critical exponent
\begin{equation}\label{eq:critical}
\qc = 1 + \frac{4s}{d+2}.
\end{equation}
We assume throughout that the potential satisfies
\begin{equation}\tag{H1}\label{H1}
V \geq 0 \quad \text{and} \quad V \in L^{d/(2s)}(\R^d),
\end{equation}
\begin{equation}\tag{H2}\label{H2}
\|V\|_{L^{d/(2s)}} < \frac{1}{8}.
\end{equation}

Our first result is global existence and decay above the critical exponent.

\begin{main}\label{thm:A}
Let $d \geq 3$ and $1/2 < s < 1$, and assume \eqref{H1} and \eqref{H2}. If $q > \qc$ and $\|u_0\|_{L^1} + \|u_0\|_{L^\infty}$ is small enough, the solution of \eqref{eq:flow} is global, and
\begin{equation}\label{eq:decay}
\|u(t)\|_{L^2} \leq C (1 + t)^{-\frac{d-2s}{4s}} \qquad \text{for all } t \geq 0,
\end{equation}
where $C$ depends only on $d$, $s$, $q$ and $u_0$.
\end{main}

Our second result is blow-up below it.

\begin{main}\label{thm:B}
Let $d \geq 3$ and $1/2 < s < 1$, and assume \eqref{H1}. If $1 < q < \qc$, every nonnegative solution of \eqref{eq:flow} with $u_0 \not\equiv 0$ blows up in finite time.
\end{main}

\Cref{thm:B} needs only \eqref{H1}: the smallness \eqref{H2} is used in \cref{prop:energy}, which only the proof of \cref{thm:A} needs. Neither theorem covers $s \leq 1/2$, where the Nash-type inequality of \cref{lem:nash} fails, or the critical exponent itself; see \cref{conj:critical}.

\section{Preliminaries}\label{sec:prelim}

We write $\|\cdot\|_r$ for the norm of $L^r(\R^d)$ and $\dot H^s$ for the homogeneous Sobolev space.

\begin{lemma}\label{lem:sobolev}
For $d > 2s$ and $f \in \dot H^s(\R^d)$, we have $\|f\|_{2d/(d-2s)} \leq S_{d,s} \|(-\Delta)^{s/2} f\|_2$.
\end{lemma}

\begin{lemma}\label{lem:hardy}
Under \eqref{H1} and \eqref{H2}, the operator $(-\Delta)^s + V$ is nonnegative, and $\langle ((-\Delta)^s + V) f, f\rangle \geq \frac{7}{8} \|(-\Delta)^{s/2} f\|_2^2$.
\end{lemma}

\section{Linear estimates}\label{sec:linear}

\begin{proposition}\label{prop:semigroup}
Assume \eqref{H1}. The semigroup generated by $-((-\Delta)^s + V)$ is positivity preserving and contracts every $L^r$, $1 \leq r \leq \infty$.
\end{proposition}

\begin{lemma}\label{lem:linear-decay}
Assume \eqref{H1} and \eqref{H2}. For $f \in L^1 \cap L^2$, the linear flow satisfies $\|e^{-t((-\Delta)^s + V)} f\|_2 \leq C\, t^{-d/(4s)} \|f\|_1$.
\end{lemma}

\begin{lemma}\label{lem:nash}
Let $1/2 < s < 1$. For $f \in L^1 \cap \dot H^s$,
\begin{equation}\label{eq:nash}
\|f\|_2^{2 + 4s/d} \leq N_{d,s}\, \|(-\Delta)^{s/2} f\|_2^2\, \|f\|_1^{4s/d}.
\end{equation}
\end{lemma}

\begin{proof}
For $s > 1/2$ the kernel of $(-\Delta)^{s}$ is integrable away from the origin, so the proof of the classical Nash inequality applies. For $s \leq 1/2$ the constant $N_{d,s}$ obtained this way is infinite.
\end{proof}

\section{Global existence and decay}\label{sec:global}

\begin{theorem}\label{thm:local}
Assume \eqref{H1}. For every $u_0 \in L^1 \cap L^\infty$ there is a unique solution of \eqref{eq:flow} on $[0, T]$, where $T \geq c_q \|u_0\|_\infty^{1-q}$.
\end{theorem}

\begin{proposition}\label{prop:energy}
Assume \eqref{H1} and \eqref{H2}, and let $q > \qc$. If $\|u_0\|_1 + \|u_0\|_\infty \leq \varepsilon_0$, then $E(t) = \|u(t)\|_2^2$ satisfies $E'(t) \leq -\frac{1}{2} N_{d,s}^{-1} E(t)^{1 + 2s/d} \|u_0\|_1^{-4s/d}$.
\end{proposition}

\begin{proof}[Proof of \cref{thm:A}]
By \cref{thm:local} the solution exists on some interval. By \cref{prop:energy} and \cref{lem:nash}, $\|u(t)\|_2$ is bounded on that interval, so the solution is global. Integrating the differential inequality of \cref{prop:energy} gives \eqref{eq:decay}.
\end{proof}

\section{Blow-up}\label{sec:blowup}

\begin{proposition}\label{prop:kaplan}
Let $\phi_R(x) = \phi(x/R)$ for a fixed nonnegative bump function $\phi$. For a nonnegative solution of \eqref{eq:flow}, $J(t) = \int u(t) \phi_R$ satisfies $J'(t) \geq -C R^{-2s} J(t) + c R^{-d(q-1)} J(t)^q$.
\end{proposition}

\begin{proof}[Proof of \cref{thm:B}]
For $q < \qc$, choosing $R$ large in \cref{prop:kaplan} makes the nonlinear term dominate, and $J$ blows up in finite time.
\end{proof}

\begin{remark}\label{rem:s-half}
At $s = 1/2$ the argument of \cref{prop:kaplan} still applies, but \cref{thm:B} is stated only for $s > 1/2$ because we use \cref{thm:local}, whose proof needs \cref{lem:nash}.
\end{remark}

At the critical exponent the two terms of \cref{prop:kaplan} balance, and our methods decide nothing.

\begin{conjecture}\label{conj:critical}
Let $d \geq 3$ and $1/2 < s < 1$, and assume \eqref{H1}. If $q = \qc$, every nonnegative solution with $u_0 \not\equiv 0$ blows up in finite time.
\end{conjecture}

\begin{thebibliography}{1}
\bibitem{Odum} K.~Odum, \emph{Anomalous diffusion with absorption}, Arch. Ration. Flow Anal. \textbf{9} (2021), 1--44.
\end{thebibliography}

\end{document}
